\SourceLeaf{14}
with the upper or lower signs applying according as $M$ lies outside or inside the body. In the former case, the portions of the cone from $P'$ to $P''$, from $P'''$ to $P''''$, and so on constitute part of the solid body. In the latter case, the portions from $M$ to $P'$, from $P''$ to $P'''$, and so on do so. In either case, then, the portion of the body lying inside the cone on base $\mathrm{d}\Sigma$ is
\[
=-\tfrac13(r'\mathrm{d}s'.\cos MQ'+r''\mathrm{d}s''.\cos MQ''+r'''\mathrm{d}s'''.\cos MQ'''+\text{ etc.})
\]
Treating all elements $\mathrm{d}\Sigma$ in the same way and summing, we obtain
\par\medskip\Needspace{4\baselineskip}
\begin{center}THEOREM FIVE.\end{center}
\par
\textit{The volume of the whole body is equal to the integral $-\tfrac13\int r\mathrm{d}s.\cos MQ$ taken over the whole surface of the body.}

\Article{8}
\par
Let us now suppose that the body is of uniform density and that each of its elements exerts on point $M$ an attraction proportional to some function of the distance. Then, with $\rho$ denoting the distance of the element from the attracted point, the attraction is expressed by the volume of the element multiplied by $f\rho$. First imagine our cone standing on base $\mathrm{d}\Sigma$, entirely filled with matter and divided into infinitely many elements by spherical surfaces infinitely close to one another and drawn about the center $M$. Such an element, at the sphere whose radius is $=\rho$, is expressed by $\rho\rho\mathrm{d}\rho.\mathrm{d}\Sigma$, and the force with which it acts on $M$ by $\mathrm{d}\Sigma.\rho\rho f\rho.\mathrm{d}\rho$. Thus, denoting the integral $\int\rho\rho f\rho.\mathrm{d}\rho$ by $\Phi\rho$, we see that $\mathrm{d}\Sigma.(\Phi\rho-\Phi0)$ expresses the attraction that the part of the cone extending from the vertex as far as distance $\rho$ exerts on point $M$, or, more generally, that $\mathrm{d}\Sigma.(\Phi\rho'-\Phi\rho)$ expresses the attraction of the cone between the distances $\rho$ and $\rho'$ from the vertex. All the parts of our body lying inside the cone will therefore attract point $M$ in direction $M\Pi$ with a force expressed by
\[
\mathrm{d}\Sigma.(-\Phi r'+\Phi r''-\Phi r'''+\text{ etc.})
\]
whenever $M$ lies outside the body, or by
\[
\mathrm{d}\Sigma.(-\Phi0+\Phi r'-\Phi r''+\Phi r'''-\text{ etc.})
\]
whenever $M$ lies inside the body, or equivalently
\par
in the former case by
